\pagestyle{fancy}
\fancyhead[R]{Física Nuclear y de Partículas: Apéndice C}
\fancyhead[L]{4º Grado en Física}
\chapter{Factor de Forma. Potencial de Yukawa}
\label{chap: Apendice C}
Vamos a considerar un scattering elástico a través de un potencial de Yukawa
\begin{equation}
    V(r)=\frac{1}{r}\exp{\dfrac{-r}{a}}
    \label{eq: potencial de yukawa}
\end{equation}
Este potencial es de gran interés en Física, como veremos más adelante. La constante $a$ es una longitud característica del sistema que en el límite de $a\xrightarrow{}\infty$ se recupera el potencial Coulombiano.\\\\
De las definiciones de la sección eficaz (diferenciamos la ecuación \ref{eq: seccion eficaz}) podemos escribir la densidad de estados finales de un proceso a dos cuerpos a través de la expresión diferencial de la Ecuación \ref{eq: densidad de estados inicial de una resonancia} 
\begin{equation}
    \dd \rho_\mathrm{f}=\frac{V}{v_i} \frac{1}{\qty(2\pi\hbar)^3}p^2 \dd \Omega
    \label{eq: diferencial de estados finales a dos cuerpos}
\end{equation}
Por lo tanto 
\begin{equation}
    \dv{\sigma}{\Omega}=\frac{V}{v_i} \frac{2\pi}{\hbar} \abs{M_\text{if}}^2 \frac{V}{v_i} \frac{1}{\qty(2\pi\hbar)^3} p^2=\frac{1}{4\pi^2 \hbar^4} m^2\gamma^2 V^2 \abs{M_\mathrm{if}}^2
\end{equation}
donde $$\abs{M_\mathrm{if}}=\displaystyle \int \frac{1}{\sqrt{V}}\exp{i\frac{\vec{k_\mathrm{f}}\cdot \vec{r}}{\hbar}}V(\vec{r}) \frac{1}{\sqrt{V}}\exp{i\frac{\vec{k_\mathrm{f}}\cdot \vec{r}}{\hbar}} \,\dd^3r=\frac{1}{V} \displaystyle \int \exp{-i\frac{\vec{q}\cdot \vec{r}}{\hbar}}V(\vec{r}) \, \dd^3r $$
donde la partícula incidente $\vec{k_\mathrm{i}}$ y la saliente $\vec{k_\mathrm{f}}$ están descritas por ondas planas (aproximación de Born). Podemos tomar $\gamma=1$, lo que no influye en la deducción, así 
\begin{equation}
    \dv{\sigma}{\Omega}=\frac{1}{4\pi^2 \hbar^4} m^2 V^2 \abs{f(\vec{q})}^2 \quad \text{con} \quad f(\vec{q})= \displaystyle \int \exp{-i\frac{\vec{q}\cdot \vec{r}}{\hbar}}V(\vec{r}) \, \dd^3r
    \label{eq: scattering de dos particulas cargadas con el propagador del potencial}
\end{equation}

Consideremos ahora que la interacción se produce con una distribución de carga. Cada elemento de volumen contiene una carga $$Z\,e\, \rho(\vec{r}) \dd^3r$$
siendo $\rho(\vec{r})$ la densida de carga normalizada. La contribución de cada elemento de volumen al potencial vendrá dada por,
\begin{equation}
    \dd V(\vec{r})=-\frac{Z\,e^2}{r^{\prime \prime}} \exp{\frac{r^{\prime \prime}}{a}} \rho(\vec{r^\prime})\, \dd^3r^\prime \quad \text{con}  \quad r^{\prime \prime}=\abs{\vec{r^\prime}}
\end{equation}
Por tanto 
\begin{equation}
    V(\vec{r})=-Ze^2\int \dd^3r^\prime \,\frac{\rho(\vec{r^\prime})}{r^{\prime \prime}} \exp{\frac{r^{\prime \prime}}{a}}
    \label{eq: Potencial con la distribucion de carga}
\end{equation}
Dado que $\vec{r}=\vec{r^\prime}+\vec{r^{\prime \prime}}$, sustituyendo en la ecuación \ref{eq: scattering de dos particulas cargadas con el propagador del potencial} obteneemos
\begin{equation}
    f(\vec{q})=-Ze^2\int \dd^3r \exp{i\frac{\vec{q}\cdot \vec{r^\prime}}{\hbar}} \rho(\vec{r^\prime}) \underbrace{\int \dd^3r \frac{1}{r^{\prime \prime}}\exp{\frac{r^{\prime \prime}}{a}}  \exp{i\frac{\vec{q}\cdot \vec{r^\prime}}{\hbar}}}_{I_2}
\end{equation}
Para un $\vec{r^\prime}$ fijo, $\dd^3r$ se puede reemplazar por $\dd^3r^{\prime \prime}$ en $I_2$ y así,
\begin{equation}
    I_2=\int^\infty_0 r^{\prime \prime } \frac{1}{r^{\prime \prime}}\exp{\frac{r^{\prime \prime}}{a}} \, \dd r^{\prime \prime}\,\underbrace{\int^\pi_0 \exp{i\frac{qr^{\prime \prime}\cos(\theta)}{\hbar}} \sin(\theta)\dd \theta }_{2\dfrac{\sin(qr^{\prime \prime}/\hbar)}{qr^{\prime \prime}/\hbar}}\,\int^{2\pi}_0 \dd \varphi
\end{equation}
Por lo tanto 
\begin{equation}
    I_2=\frac{4\pi}{q/\hbar} \int^\infty_0 \exp{\frac{r^{\prime \prime}}{a}} \sin\qty(\frac{qr^{\prime \prime}}{\hbar}) \dd r^{\prime \prime}=\frac{4\pi}{\qty(\dfrac{q}{\hbar})^2+\qty(\dfrac{1}{a})^2}
    \label{eq: transformada de fourier potencial de yukawa}
\end{equation}
Donde hemos utilizado la integral de Bessel para hayar el resultado final.
Si tomamos el límite de $a\xrightarrow[]{} \infty \Rightarrow I_2=\dfrac{4\pi \hbar^2}{q^2}$ 
Por tanto
\begin{equation}
    f(\vec{q})=\underbrace{Ze^2\dfrac{4\pi \hbar^2}{q^2}}_{V(q)} \underbrace{\int \dd^3r \exp{i\frac{\vec{q}\cdot \vec{r}}{\hbar}} \rho(\vec{r})}_{F(\vec{q})}
\end{equation}
Finalmente
\begin{equation}
    \dv{\sigma}{\Omega}=\frac{1}{4\pi^2\hbar^4} m^2 \abs{V(q)}^2\abs{F(\vec{q})}^2
    \label{eq: Dispersion EM a dos cuerpos con factor de forma}
\end{equation}
Si la distribución fuese puntual $\rho(\vec{r})=\delta(\vec{r})$ y $F(\vec{q})=1$. Sustituyendo el potencial en la ecuación \ref{eq: Dispersion EM a dos cuerpos con factor de forma} obtenemos
\begin{equation}
      \dv{\sigma}{\Omega}=4Z^2 \frac{m^2e^4}{q^4} \abs{F(\vec{q})}^2
    \label{eq: Dispersion relativista EM a dos cuerpos con factor de forma}
\end{equation}
\clearpage
\subsubsection{Información del factor de forma. Radio cuadrático medio}
Estudiemos un poco más en detalle la información que contiene el factor de forma. Para ello utilizaremos 
\begin{equation}
    \exp{i\vec{q}\cdot \vec{r}}= 4\pi \sum^\infty_{\ell=0} \sum^\ell_{m=-\ell} i^\ell j_\ell(q,r)Y_{\ell m}(\Omega_q)Y_{\ell m}(\Omega_r)
    \label{eq: desarrollo en armonicos esfericos exponencial imaginaria}
\end{equation}
Suponiendo que la distribución de carga es esférica $\rho(\vec{r})=\rho(r)$
\begin{equation}
    F(q)= 4\pi \sum^\infty_{\ell=0} \sum^\ell_{m=-\ell} i^\ell \int^\infty_0 r^2 j_\ell(qr) \rho(r)Y_{\ell m}(\Omega_q) \dd r\underbrace{\int Y_{\ell m}(\Omega_r)\dd\Omega_r}_{\sqrt{4\pi}\delta_{\ell 0} \delta_{m 0}}
\end{equation}
Las deltas eliminan los sumatorios y la unidad imaginaria, además como $\ell=0,m=0$, obtenemos
\begin{equation}
    F(q)=4\pi\int^\infty_0 r^2 j_0(qr)\rho(r)\dd r=\frac{4\pi}{q} \int^\infty_0 r \sin(qr)\rho(r)\dd r
    \label{eq: factor de forma deistribucion esferica}
\end{equation}
Estudiemos el límite en el que $q\xrightarrow{} 0 \Rightarrow k_\mathrm{i}\approx k_\mathrm{f}$, es decir, la partícula incidente interacciona muy poco con la distribución de carga. Como $r$ varía sobre el intervalo finito de la distribución
$$q\xrightarrow[]{}0 \Rightarrow qr\xrightarrow{}0 \Rightarrow \sin(qr)|_{q\xrightarrow[]{}0}=qr-\frac{1}{6}q^3r^3+\cdots$$
Así
\begin{equation*}
    \begin{aligned}
        F(q)=&\frac{4\pi}{q}\int^\infty_0 r\qty[qr-\frac{1}{6}q^3r^3+\cdots]\rho(r)\dd r=\frac{4\pi}{q}\int^\infty_0 q\rho(r) r^2\dd r -\frac{1}{6}\frac{4\pi}{q}\int^\infty_0 q^3\rho(r)r^4 \dd r+ \cdots\\
        &=4\pi\int^\infty_0 q\rho(r) r^2\dd r -\frac{1}{6}4\pi q^2\int^\infty_0 \rho(r)r^4 \dd r+ \cdots=\int \rho(r)\dd^3r-\frac{1}{6}q^2\int r^2 \rho(r)\dd^3r
    \end{aligned}
\end{equation*}
Como la distribución está normalizada y la última integral es $\expval{r^2}$ por lo tanto 
\begin{equation}
    F(q)|_{q\xrightarrow{}0}=1-\frac{1}{6} q^2 \expval{r^2} \xrightarrow{q\xrightarrow{}0}1
\end{equation}
Por lo que podemos hayar el radio cuadrático medio de la distribución con el factor de forma
\begin{equation}
    \expval{r^2}=-6\left.\dv{F(q)}{q^2}\right |_{q\xrightarrow{}0}
\end{equation}